\begin{center}
\hypertarget{native-038-article021}{}\pdfbookmark[2]{21.}{article-021}21.
\end{center}
In HYPERBOLA $\varphi$ atque $E$ quantitates imaginariae fierent, quales si aversamur, illarum loco aliae quantitates auxiliares sunt introducendae. Angulum cuius cosinus $=\frac{1}{e}$ iam supra per $\psi$ designavimus, radiumque vectorem
\[
=\frac{p}{2e\cos\frac{1}{2}(v-\psi)\cos\frac{1}{2}(v+\psi)}
\]
invenimus. Factores in denominatore huius fractionis, $\cos\frac{1}{2}(v-\psi)$ et $\cos\frac{1}{2}(v+\psi)$, aequales fiunt pro $v=0$, secundus evanescit pro valore maximo positivo ipsius $v$, primus vero pro valore maximo negativo. Statuendo igitur $\frac{\cos\frac{1}{2}(v-\psi)}{\cos\frac{1}{2}(v+\psi)}=u$, erit $u=1$ in perihelio; crescet in infinitum, dum $v$ ad limitem suum $180^0-\psi$ appropinquat; contra decrescet in infinitum, dum $v$ ad limitem \SourceLeaf{039}alterum $-(180^0-\psi)$ regredi supponitur: quod fiet ita, ut valoribus oppositis ipsius $v$ valores reciproci ipsius $u$, vel quod idem est valores tales, quorum logarithmi oppositi sunt, respondeant.

Hic quotiens $u$ percommode in hyperbola ut quantitas auxiliaris adhibetur; aequali fere concinnitate istius vice fungi potest angulus cuius tangens $=\operatorname{tang}\frac{1}{2}v\, .\,\BareRoot{\frac{e-1}{e+1}}$, quem ut analogiam cum ellipsi sequamur, per $\frac{1}{2}F$ denotabimus. Hoc modo facile sequentes relationes inter quantitates $v$, $r$, $u$, $F$ colliguntur, ubi $a=-b$ statuimus, ita ut $b$ evadat quantitas positiva.
\[
\mathrm{I.}\quad b=p\operatorname{cotang}\psi^2
\]
\[
\mathrm{II.}\quad r=\frac{p}{1+e\cos v}=\frac{p\cos\psi}{2\cos\frac{1}{2}(v-\psi)\cos\frac{1}{2}(v+\psi)}
\]
\[
\begin{aligned}
\mathrm{III.}\quad\operatorname{tang}\frac{1}{2}F&=\operatorname{tang}\frac{1}{2}v\, .\,\BareRoot{\frac{e-1}{e+1}}\\
&=\operatorname{tang}\frac{1}{2}v\operatorname{tang}\frac{1}{2}\psi=\frac{u-1}{u+1}
\end{aligned}
\]
\[
\begin{aligned}
\mathrm{IV.}\quad u&=\frac{\cos\frac{1}{2}(v-\psi)}{\cos\frac{1}{2}(v+\psi)}=\frac{1+\operatorname{tang}\frac{1}{2}F}{1-\operatorname{tang}\frac{1}{2}F}\\
&=\operatorname{tang}(45^0+\frac{1}{2}F)
\end{aligned}
\]
\[
\begin{aligned}
\mathrm{V.}\quad\frac{1}{\cos F}&=\frac{1}{2}\left(u+\frac{1}{u}\right)\\
&=\frac{1+\cos\psi\cos v}{2\cos\frac{1}{2}(v-\psi)\cos\frac{1}{2}(v+\psi)}=\frac{e+\cos v}{1+e\cos v}\cdot
\end{aligned}
\]

Subtrahendo ab aeq. V utrimque $1$, prodit
\[
\begin{aligned}
\mathrm{VI.}\quad\sin\frac{1}{2}v\, .\,\BareRoot{r}&=\sin\frac{1}{2}F\, .\,\BareRoot{\frac{p}{(e-1)\cos F}}\\
&=\sin\frac{1}{2}F\, .\,\BareRoot{\frac{(e+1)b}{\cos F}}\\
&=\frac{1}{2}(u-1)\BareRoot{\frac{p}{(e-1)u}}=\frac{1}{2}(u-1)\BareRoot{\frac{(e+1)b}{u}}\cdot
\end{aligned}
\]

Simili modo addendo utrimque $1$ fit
\[
\begin{aligned}
\mathrm{VII.}\quad\cos\frac{1}{2}v\, .\,\BareRoot{r}&=\cos\frac{1}{2}F\, .\,\BareRoot{\frac{p}{(e+1)\cos F}}\\
&=\cos\frac{1}{2}F\, .\,\BareRoot{\frac{(e-1)b}{\cos F}}\\
&=\frac{1}{2}(u+1)\BareRoot{\frac{p}{(e+1)u}}=\frac{1}{2}(u+1)\BareRoot{\frac{(e-1)b}{u}}\cdot
\end{aligned}
\]

Dividendo VI per VII ad III reveniremus; multiplicatio producit
\[
\begin{aligned}
\mathrm{VIII.}\quad r\sin v&=p\operatorname{cotang}\psi\operatorname{tang}F=b\operatorname{tang}\psi\operatorname{tang}F\\
&=\frac{1}{2}p\operatorname{cotang}\psi\, .\,\left(u-\frac{1}{u}\right)=\frac{1}{2}b\operatorname{tang}\psi\, .\,\left(u-\frac{1}{u}\right).
\end{aligned}
\]

E combinatione aequatt. II, V porro facile deducitur
\[
\mathrm{IX.}\quad r\cos v=b\left(e-\frac{1}{\cos F}\right)=\frac{1}{2}b\left(2e-u-\frac{1}{u}\right)
\]
\[
\mathrm{X.}\quad r=b\left(\frac{e}{\cos F}-1\right)=\frac{1}{2}b\left\{e\left(u+\frac{1}{u}\right)-2\right\}.
\]
